\uuid{5j65}
\exo7id{51}
\titre{exo7 51}
\auteur{bodin}
\organisation{exo7}
\datecreate{1998-09-01}
\isIndication{false}
\isCorrection{false}
\chapitre{Nombres complexes}
\sousChapitre{Racine n-ieme}
\cpge{mpsi-mp2i.nombres-complexes.f, pcsi.nombres-complexes.f, ptsi.nombres-complexes.f, mpsi-mp2i.calcul-algebrique-trigonometrie.a, pcsi.calcul-algebrique-trigonometrie.a, ptsi.calcul-algebrique-trigonometrie.a, bcpst1.methodes-calcul}
\module{Algèbre}
\niveau{L1}
\difficulte{}

\contenu{
\question{
Soit $\epsilon$ une racine $n$-i\`eme de l'unit\'e ; calculer
$$ S = 1 +2\epsilon+3\epsilon^2+\cdots+n\epsilon^{n-1}.$$
}
}
